\chapter{LITERATURE REVIEW}

Hand-drawn circuit recognition and digitization, while not a vastly researched area, has a few notable examples of its implementation. Most of the study and implementation related to this has been done on classifying the different circuit components (like battery, resistors, inductors, etc.) and there's not much on arguably the more difficult aspect of circuit recognition: node and wire connectivity. For instance, in a 2021 study, 20 different components were classified from self-made hand-drawn circuits with the recognition accuracy of 97.33\% \cite{adnan2021cnnbased}.


In one of the oldest studies on the automatic recognition of hand-drawn circuits, an approach of identifying the components first and then the node segments (taking into consideration their spatial relationship) was implemented, achieving an accuracy of 86\% for components and 92.5\% for the nodes. The node segments were then traced assuming wire between them, considering a variety of hard-coded conditions \cite{edwards2000machine}. In a recent 2024 research, a couple of \gls{uk} engineers were successful in digitizing images of electrical circuit schematics using a novel pattern-recognition methodology and custom-trained \gls{ocr} model achieving respective success rates of 86.4\% and 99.6\% \cite{kelly2024digitizing}. This, however, was trained on printed/computer-made circuit diagrams, not hand-drawn ones and also isn't able to understand the values of the circuit components.



To overcome the limitations of \glspl{cnn} in understanding circuit topology, graph-based methods have gained traction. Circuit2Graph, a pipeline that transforms circuit diagrams into graph structures where nodes represent components and edges denote electrical connections was introduced in 2023. Using \glspl{gnn}, their system achieved 97.9\% accuracy on circuit classification tasks, outperforming traditional methods \cite{yamakaji2023circuit2graph}. This approach demonstrated that \glspl{gnn} can effectively model the global structure of a circuit, even when visual features alone are insufficient. Similarly, Bayer et al.\ \cite{bayer2023instance} used Mask \gls{rcnn} for instance segmentation of components in hand-drawn diagrams, followed by graph extraction techniques to infer connections. Their model achieved around 94\% mask accuracy, but they noted that symbol recognition was ``reasonable, yet incomplete,'' especially in more complex drawings or where handwriting was less consistent. A paper published in the International Research Journal of Engineering and Technology proposed a two-step method for electronic circuit diagram detection and recognition; first detecting and classifying the components present in the circuit and then labeling the classified components to form a string which is then fed to an \gls{fsm} system for circuit recognition \cite{abhinav2020electric}.



With the use of YOLOv5 for detection of circuit components and Probabilistic Hough transform-based solution for node recognition and taking in hand-drawn circuits, a couple of Indian engineers were successful in achieving not only a \gls{map} of 98.2\% in component detection but also a near-real time generation of circuit schematic with an accuracy of 80\% \cite{reddy2021handdrawn}. However, their circuits were not labelled which is to say the values for the battery (in volts), resistance (in Ohms), etc.\ were not written in the circuit. Another very recent study suggested a three-step process to detect and identify the circuit and its components with the steps including detection of components, detection of lines and their junctions and then the detection of text placed near the components and then combining the data from these three streams into a \gls{json} format to eventually derive a netlist from \cite{hemker2024schematics}. This is a similar workflow to the one this project intends to use; however, the scope of our project is much larger owing to the fact that an editable schematic is to be delivered on top of taking in hand-written circuit diagrams as opposed to generating a netlist from a digital schematic.

Netlist generation is also a very critical aspect of this domain. While traditional \gls{eda} tools such as KiCad and \gls{ltspice} support netlist export from manually created schematics, they do not allow direct image-to-netlist conversion \cite{kicad2024spice}. This limitation was addressed with Img2Sim; an artificial neural network-based system that detects components, infers connectivity, and outputs a complete \gls{spice} netlist. Their model reported over 90\% accuracy in topology extraction and represents one of the first functional pipelines to translate a hand-drawn circuit into a ready-to-simulate format \cite{sertdemir2022img2sim}. However, Img2Sim and similar systems typically require clean input images and are limited to a predefined set of components, reducing their robustness in real-world use cases. While the feasibility of such pipelines is now evident, a generalized, accurate, and open-source solution that goes the extra mile to also visualize the circuits digitally remains unavailable.



Parallel to recognition and netlist generation, visual question answering (\gls{vqa}) for circuits is emerging as a useful educational application. CircuitVQA is a dataset containing more than 115,000 question-image pairs based on 5,725 unique circuits. It covers approximately 70 symbol types and supports the training of transformer-based \gls{vqa} models capable of answering questions like ``What is the function of R1?'' or ``Is this a voltage divider?'' Early experiments demonstrated that models fine-tuned on CircuitVQA outperformed generic \gls{vqa} systems, indicating the importance of domain-specific training. Meanwhile, large language models (\glspl{llm}) such as GPT-4 have also been explored for circuit explanation tasks. For instance, the SPICEPilot framework attempts to generate and interpret \gls{spice} code using GPT-based models \cite{vungarala2024spicepilot}. However, these general-purpose models often fail to respect domain conventions, such as valid W/L ratios or simulation types, resulting in inaccurate outputs.



Finally, simulation integration plays a vital role in validating and interacting with generated circuits. Ngspice, an open-source simulation engine, is widely used in research and is natively supported in KiCad. It enables \gls{dc}, \gls{ac}, and transient analysis, making it suitable for most educational and prototyping needs. In contrast, commercial tools like Simulink and PSpice provide broader capabilities but are cost-prohibitive. Some research efforts, including Img2Sim, have leveraged open standards to generate \gls{spice} files compatible with tools like Ngspice or PSpice. Other projects utilize Python-based libraries such as PySpice to enable automated simulation directly within a code-driven environment which is exactly how this project intends to accomplish its simulation needs.



While significant progress has been made in component recognition and value comprehension, connectivity inference, netlist generation, and interactive Q\&A, existing systems remain fragmented or limited in scope. An end-to-end system that combines all these elements into a user-friendly platform remains a critical need for both educational and practical applications.